%%%%%%%%%%%%%%%%%%%%%%%%%%%%%%%%%%%%%%%%%%%%%%%%%%%%%%%%%%%%%%%%%%%%%%%%%%%
% Englische Fassung des Uebungsblatts zu Lehreinheit 8 (Teil A und B).
% Vorlage: handouts/LE08.tex. Nur in handouts.tex mit \def\wuzsprache{en}.
%%%%%%%%%%%%%%%%%%%%%%%%%%%%%%%%%%%%%%%%%%%%%%%%%%%%%%%%%%%%%%%%%%%%%%%%%%%
\begin{uebungsblatt}{Cognitive Biases}

\aufgabe{Warm-up: a fictitious short report}[in class, individual work, 5 min.]
\textit{The following report is fictitious and was written for this
exercise.}

\begin{center}
\fbox{\parbox{0.92\textwidth}{\small
\textbf{\enquote{Baroque music boosts concentration}}\\[1mm]
(1)~A lecturer who has offered concentration training with music for
years wants to show that baroque music improves concentration. (2)~She
recruits 24 students from her own seminar. (3)~All of them first solve
tasks without music in the morning and then tasks with music in the
afternoon of the same day. (4)~The lecturer marks the solutions herself;
each sheet states whether it was produced with or without music.
(5)~Twelve measures of concentration are recorded. One of them shows a
significant advantage for the music condition; the other eleven are not
reported. (6)~Because the result after 24 people was \enquote{just short of
significance}, ten more people were tested. (7)~The discussion states that
the hypothesis was clear from the outset and confirms the author's
many years of practical experience. (8)~An earlier study by the author
that found no effect was never submitted.
}}
\end{center}

For each number, note whether a bias (a systematic distortion in a
particular direction) could be at work there, which one, and in which
phase of the research process: planning, data collection, analysis,
interpretation or publication.
\schreiblinien{5}

\aufgabe{Bias hunt in an article}[in class, pairs, 25 min.; finish as homework]
Article (chosen by the lecturer): \hrulefill

Find at least two places where a bias could have influenced the result.
Check, one after the other: Who conducted the study, and with what
interest? How were participants recruited and assigned? Who measured, and
did that person know the condition? How many measures and analysis
variants were there, and are all of them reported? Was there a
preregistration? Does the interpretation go beyond the data?

\begin{center}
\small
\renewcommand{\arraystretch}{2.6}
\begin{tabular}{|p{0.13\textwidth}|p{0.11\textwidth}|p{0.15\textwidth}|p{0.27\textwidth}|p{0.18\textwidth}|}
\hline
\textbf{Page / quote} & \textbf{Phase} & \textbf{Bias} & \textbf{Reason: in which direction does it shift the result?} & \textbf{Remedy (used?)} \\
\hline
 & & & & \\ \hline
 & & & & \\ \hline
 & & & & \\ \hline
\end{tabular}
\end{center}

What did the authors do well methodologically?
\schreiblinien{2}
\end{uebungsblatt}

\begin{uebungsblattb}{3}{Bias in Your Own Project}

\aufgabe{Reflection: the bias in my own research process}[homework]
Base this on your design sketch from Unit~7.

\begin{center}
\renewcommand{\arraystretch}{1.8}
\begin{tabular}{|p{0.36\textwidth}|p{0.54\textwidth}|}
\hline
Which bias is the greatest threat? & \\ \hline
In which phase? & \\ \hline
How would I notice that it has had an effect? & \\ \hline
With which \emph{procedure} do I limit it? & \\ \hline
\end{tabular}
\end{center}

\aufgabe{Diary of wonder}[homework]
Describe a situation in which you knew something and still perceived or
decided differently. What would have helped you follow your knowledge
rather than your experience?

\begin{lehrhinweis}[Expected answers]
\textbf{Task~1.} (1)~Conflict of interest, confirmation bias in planning.
(2)~Selection and dependence: her own students, possible effects of
wanting to please. (3)~The condition is confounded with order and time of
day (practice, fatigue); not a bias in the narrow sense, but a threat to
internal validity that can work in the direction of the hoped-for effect.
(4)~Experimenter effect, no blind marking (data collection).
(5)~Selective reporting, p-hacking across twelve measures (analysis).
(6)~Optional stopping, collecting more data until significance
(analysis). (7)~Hindsight bias or HARKing, no preregistration
(interpretation). (8)~Publication bias (publication).
\textbf{Task~2.} At least two places with location, named distortion,
justified direction and a concrete remedy. Mere weaknesses (small sample,
convenience sample) count as uncertainty or a limit to external validity;
they count as bias only if a direction is justified.
\textbf{Task~3.} A procedure is required (blind marking, preregistration,
a second coder, a stopping criterion fixed in advance), not a good
intention.
\end{lehrhinweis}
\end{uebungsblattb}
